%; whizzy-master slides.tex
%; whizzy frame -pdf ./skimopen 

\usepackage{xspace} % To get the right spacings in front of : and so on
\usepackage[english]{babel}
\usepackage{subfigure}
\usepackage{bussproofs}
\usepackage{pgf}
\usepackage{colortbl}
\usepackage{xcolor}
\usepackage{abbrevs}
\usepackage[color]{coqdoc}
\def\coqlibrary#1#2{}

\usepackage{hyperref}
\usepackage{coq}
\usepackage{code}
\usepackage{codecolor}
\usepackage{cond}
\usepackage{me}
\usepackage{tikz}
\usetikzlibrary{shapes}
\usepackage{pifont}
\usepackage{proof}
\usepackage{utf}
\usepackage{natbib}
\usepackage[OT1]{fontenc}
\usepackage{alltt}
\usepackage{amsthm}


%%% MMCG symbols

\newcommand{\idot}[1]{\mbox{$\bullet_{#1}$}}
\newcommand{\baredot}{\mbox{$\bullet$}}
\newcommand{\iback}[1]{\mbox{$\backslash_{#1}$}}
\newcommand{\back}{\mbox{$\backslash$}}
\newcommand{\islash}[1]{\mbox{$/_{#1}$}}
\newcommand{\icomma}[3]{\mbox{$(#1\,,\,#2)_{#3}$}}
%\newcommand{\icomma}[3]{\mbox{$#1\,\mathbin{\circ}_{#3}\,#2$}}
%\newcommand{\comma}{\mbox{$\mathbin{\circ}$}}
\newcommand{\comma}{\mbox{,}}

\newcommand{\slashe}{\mbox{\textbf{$/_{\textbf{E}}$}}}
\newcommand{\slashi}{\mbox{\textbf{$/_{\textbf{I}}$}}}
\newcommand{\backe}{\mbox{\textbf{$\backslash_{\textbf{E}}$}}}
\newcommand{\backi}{\mbox{\textbf{$\backslash_{\textbf{I}}$}}}

\newcommand{\dote}{\mbox{\textbf{$\baredot_{\textbf{E}}$}}}
\newcommand{\doti}{\mbox{\textbf{$\baredot_{\textbf{I}}$}}}
\newcommand{\diame}{\mbox{\textbf{$\Diamond_{\textbf{E}}$}}}
\newcommand{\diami}{\mbox{\textbf{$\Diamond_{\textbf{I}}$}}}
\newcommand{\struct}{\mbox{\textbf{struct}}}
\newcommand{\axrule}{\mbox{\textbf{Ax}}}
\newcommand{\boxe}{\mbox{\textbf{$\Box_{\textbf{E}}$}}}
\newcommand{\boxi}{\mbox{\textbf{$\Box_{\textbf{I}}$}}}
\newcommand{\AssD}{\textbf{L$_\Diamond$}}
\newcommand{\ComD}{\textbf{P$_\Diamond$}}

\newcommand{\marginok}[1]{\marginpar{\raggedright OK:#1}}
\newcommand{\tab}{{\null\hskip1cm}}
%\newcommand{\Ltac}{\mbox{\emph{$\cal L$}tac}}
\newcommand{\coq}{\mbox{\emph{Coq}}}
\newcommand{\pcoq}{\mbox{\emph{Pcoq}}}
\newcommand{\grail}{\mbox{\emph{Grail}}}
\newcommand{\lcf}{\mbox{\emph{LCF}}}
\newcommand{\hol}{\mbox{\emph{HOL}}}
\newcommand{\pvs}{\mbox{\emph{PVS}}}
\newcommand{\icharate}{\mbox{\emph{Icharate}}}
\newcommand{\isabelle}{\mbox{\emph{Isabelle}}}
\newcommand{\prolog}{\mbox{\emph{Prolog}}}
\newcommand{\goalbar}{\tt{}{\color{black}------------------------------------}\it}
\newcommand{\gallina}{\mbox{\emph{Gallina}}}
\newcommand{\joker}{\texttt{\_}}
\newcommand{\eprime}{\(\e^{\prime}\)}
\newcommand{\Ztype}{\citecoq{Z}}
\newcommand{\propsort}{\citecoq{Prop}}
\newcommand{\setsort}{\citecoq{Set}}
\newcommand{\typesort}{\citecoq{Type}}
\newcommand{\ocaml}{\mbox{\emph{OCAML}}}
\newcommand{\haskell}{\mbox{\emph{Haskell}}}
\newcommand{\why}{\mbox{\emph{Why}}}
\newcommand{\Pascal}{\mbox{\emph{Pascal}}}

\newcommand{\ml}{\mbox{\emph{ML}}}

\newcommand{\scheme}{\mbox{\emph{Scheme}}}
\newcommand{\lisp}{\mbox{\emph{Lisp}}}

\newcommand{\implarrow}{\mbox{$\Rightarrow$}}
\newcommand{\metavar}[1]{?#1}
\newcommand{\notincoq}[1]{#1}
\newcommand{\coqscope}[1]{\%#1}
\newcommand{\arrow}{\mbox{$\rightarrow$}}
\newcommand{\fleche}{\arrow}
\newcommand{\funarrow}{\mbox{$\Rightarrow$}}
\newcommand{\ltacarrow}{\funarrow}
% \newcommand{\coqand}{\mbox{\(\wedge\)}}
% \newcommand{\coqor}{\mbox{\(\vee\)}}
% \newcommand{\coqnot}{\mbox{\(\neg\)}}
\newcommand{\coqand}{\texttt{/\textbackslash}}
%\newcommand{\coqor}{\texttt{\textbackslash/}}
\newcommand{\coqnot}{\mbox{\(\neg\)}}
\newcommand{\hide}[1]{}
\newcommand{\hidedots}[1]{...}
\newcommand{\sig}[3]{\texttt{\{}#1\texttt{:}#2 \texttt{|} #3\texttt{\}}}
\renewcommand{\neg}{\mbox{$\sim$}}

%%% Operateurs, etc.
\newcommand{\impl}{\mbox{$\rightarrow$}}
\newcommand{\appli}[2]{\mbox{\tt{#1 #2}}}
\newcommand{\applis}[1]{\mbox{\texttt{#1}}}
\newcommand{\abst}[3]{\mbox{\tt{fun #1:#2 \funarrow #3}}}
\newcommand{\coqle}{\mbox{$\leq$}}
\newcommand{\coqge}{\mbox{$\geq$}}
\newcommand{\coqdiff}{\mbox{$\neq$}}
\newcommand{\coqiff}{\mbox{$\leftrightarrow$}}
\newcommand{\prodsym}{\mbox{\(\forall\,\)}}
\newcommand{\exsym}{\mbox{\(\exists\,\)}}
\newcommand{\included}{\mbox{\(\subseteq\)}}

\newcommand{\substsign}{/}
\newcommand{\subst}[3]{\mbox{#1\{#2\substsign{}#3\}}}
\newcommand{\anoabst}[2]{\mbox{\tt[#1]#2}}
\newcommand{\letin}[3]{\mbox{\tt let #1:=#2 in #3}}
\newcommand{\prodep}[3]{\mbox{\tt \(\forall\,\)#1:#2,$\,$#3}}
\newcommand{\prodplus}[2]{\mbox{\tt\(\forall\,\)$\,$#1,$\,$#2}}
\renewcommand{\dom}[1]{\textrm{dom}(#1)} % domaine d'un contexte (log function)
\newcommand{\norm}[1]{\textrm{n}(#1)} % forme normale (log function)
\newcommand{\coqZ}[1]{\mbox{\tt{`#1`}}}
\newcommand{\coqnat}[1]{\mbox{\tt{#1}}}
\newcommand{\coqcart}[2]{\mbox{\tt{#1*#2}}}
\newcommand{\alphacong}{\mbox{$\,\cong_{\alpha}\,$}} % alpha-congruence
\newcommand{\betareduc}{\mbox{$\,\rightsquigarrow_{\!\beta}$}\,} % beta reduction
%\newcommand{\betastar}{\mbox{$\,\Rightarrow_{\!\beta}^{*}\,$}} % beta reduction
\newcommand{\deltareduc}{\mbox{$\,\rightsquigarrow_{\!\delta}$}\,} % delta reduction
\newcommand{\dbreduc}{\mbox{$\,\rightsquigarrow_{\!\delta\beta}$}\,} % delta,beta reduction
\newcommand{\ireduc}{\mbox{$\,\rightsquigarrow_{\!\iota}$}\,} % delta,beta reduction

\newcommand{\slam}[2] {\mbox{$\lambda_#1\;#2$}}
\newcommand{\srho}[2] {\mbox{$\rho_#1\;#2$}}
\newcommand{\Vforall}[2] {\mbox{$\forall\,#1\,\in\,V(#2)$}}
\newcommand{\Eforall}[2] {\mbox{$\forall\,#1\,\in\,E(#2)$}}
\newcommand{\Vforone}[2] {\mbox{$\exists!\,#1\,\in\,V(#2)$}}
\newcommand{\Eforone}[2] {\mbox{$\exists!\,#1\,\in\,E(#2)$}}
\newcommand{\Vforsome}[2] {\mbox{$\exists\,#1\,\in\,V(#2)$}}
\newcommand{\Eforsome}[2] {\mbox{$\exists\,#1\,\in\,E(#2)$}}


% jugement de typage
\newcommand{\these}{\mbox{$\boldsymbol{\large \vdash}$}}
\newcommand{\rthese}[1]{\mbox{$\boldsymbol{\vdash_{#1}}$}}
\newcommand{\replace}[2]{\mbox{$#1[#2]$}}
\newcommand{\msubst}[2]{\mbox{$#1\{#2\}$}}
\newcommand{\disj}{\mbox{$\backslash/$}}
\newcommand{\conj}{\mbox{$/\backslash$}}
\newcommand{\deriv}[2]{\mbox{$#1\;\these\;#2$}}
\newcommand{\smalljuge}[3]{\mbox{$#1 \these #2 \boldsymbol{:} #3 $}}
%\newcommand{\juge}[3]{\mbox{$#1 \these #2 \boldsymbol{:} #3 $}}
\newcommand{\goal}[3]{\mbox{$#1,#2 \these^{\!\!\!?} #3  $}}
\newcommand{\sgoal}[2]{\mbox{$#1\,\these^{\!\!\!\!?}\, #2 $}}
\newcommand{\sequent}[2]{\mbox{$#1\;\these\; #2 $}}

\newcommand{\reduc}[5]{\mbox{$#1,#2 \these #3 \rhd_{#4}#5 $}}
\newcommand{\convert}[5]{\mbox{$#1,#2 \these #3 =_{#4}#5 $}}
\newcommand{\convorder}[5]{\mbox{$#1,#2 \these #3\leq _{#4}#5 $}}
\newcommand{\wouff}[2]{\mbox{$\emph{WF}(#1)[#2]$}}




\newcommand{\type}{\boldsymbol{:}}

% jugement absolu

%\newcommand{\ajuge}[2]{\mbox{$ \boldsymbol{\vdash} #1 : #2 $}}
\newcommand{\ajuge}[2]{\mbox{$\these #1 \boldsymbol{:} #2 $}}

%%% logique minimale
\newcommand{\propzero}{\mbox{$P_0$}} % types de Fzero

%%% logique propositionnelle classique
\newcommand {\ff}{\boldsymbol{f}} % faux
\newcommand {\vv}{\boldsymbol{t}} % vrai

\newcommand{\verite}{\mbox{$\cal{B}$}} % {\ff,\vv}
\newcommand{\sequ}[2]{\mbox{$#1 \vdash #2 $}} % sequent
\newcommand{\strip}[1]{#1^o} % enlever les variables d'un contexte



%%% tactiques
\newcommand{\decomp}{\delta} % decomposition
\newcommand{\recomp}{\rho} % recomposition

%%% divers
\newcommand{\cqfd}{\mbox{\textbf{cqfd}}}
\newcommand{\fail}{\mbox{\textbf{F}}}
\newcommand{\succes}{\mbox{$\blacksquare$}}
%%% Environnements


%% Fzero
\newcommand{\con}{\mbox{$\cal C$}}
\newcommand{\var}{\mbox{$\cal V$}}

\newcommand{\atomzero}{\mbox{${\cal A}_0$}} % types de base de Fzero
\newcommand{\typezero}{\mbox{${\cal T}_0$}} % types de Fzero
\newcommand{\termzero}{\mbox{$\Lambda_0$}} % termes de Fzero 
\newcommand{\conzero}{\mbox{$\cal C_0$}} % contextes de Fzero 

\newcommand{\buts}{\mbox{$\cal B$}} % buts

%%% for drawing terms
% abstraction [x:t]e
\newcommand{\PicAbst}[3]{\begin{bundle}{\bf abst}\chunk{#1}\chunk{#2}\chunk{#3}%
                        \end{bundle}}

% the same in DeBruijn form
\newcommand{\PicDbj}[2]{\begin{bundle}{\bf abst}\chunk{#1}\chunk{#2}
                       \end{bundle}}


% applications
\newcommand{\PicAppl}[2]{\begin{bundle}{\bf appl}\chunk{#1}\chunk{#2}%
                         \end{bundle}}

% variables
\newcommand{\PicVar}[1]{\begin{bundle}{\bf var}\chunk{#1}
                          \end{bundle}}

% constantes
\newcommand{\PicCon}[1]{\begin{bundle}{\bf const}\chunk{#1}\end{bundle}}

% arrows
\newcommand{\PicImpl}[2]{\begin{bundle}{\impl}\chunk{#1}\chunk{#2}%
                         \end{bundle}}



%%%% scripts coq
\newcommand{\prompt}{\mbox{\sl Coq $<\;$}}
\newcommand{\natquicksort}{\texttt{nat\_quicksort}}
\newcommand{\citecoq}[1]{\mbox{\texttt{#1}}}
\newcommand{\safeit}{\it}
\newcommand{\byrule}[1]{\mbox{$\underset{\scriptstyle #1}{\longrightarrow}$}}
\newcommand{\byrules}[1]{\mbox{$\underset{\scriptstyle #1}{\longrightarrow^*}$}}


\newcommand{\backback}{\mbox{$\backslash\backslash$}}

\renewcommand{\cons}{\mbox{::}}
%\input{univmacros}

\def\vec#1{\overrightarrow{#1}}
\def\thetitle{\Coq 8.5: what's new, what's next?}
\def\thetextitleabbr{\thetitle}
\def\thetextitle{\thetextitleabbr}
\def\thesubject{Theoretical Computer Science}
\def\thekeywords{Coq}
\def\theevent{\Coq POPL Meeting\\
January 21st 2014\\
San Diego, USA}
\def\theeventabbr{}
\def\theauthorabbr{M. Dénès \& M. Sozeau}
\usepackage{prelude-beamer}

\author{M. Dénès, \myname{}}
\def\theauthor{M. Dénès$^{\dagger}$ \& M. Sozeau$^{*}$}
\def\PI{\name{PI}}
\def\shaded#1{{\color{black!50}#1}}
 
\def\seq{\fCenter}

\def\lang{en}
\input{mathenv}
\setboolean{displayLabels}{true}

\pgfdeclareimage[height=1.4cm]{inria-logo}{inria-logo}

\pgfdeclareimage[height=1.4cm]{coq-biglogo}{barron-logo-pdf}
\pgfdeclareimage[height=0.4cm]{coq-smalllogo}{barron-logo}
\pgfdeclareimage[height=1cm,interpolate=true]{coq-logo}{barron-logo-pdf}

\pgfdeclareimage[height=7cm]{coq-share}{coq-share}

\def\imgheight{7.5cm}

\def\typea{\vdash}
\def\suba{\rightslice}
\def\indent{\coqdocindent}
\def\var{\coqdocvar}

\def\coqdocid#1{\coqdocvar{#1}}

%\renewcommand{\Type}[1]{\kw{Type}_{#1}}

\newcommand{\elt}[4]{\ensuremath{\constr{exist}_{#1,#2}~#3~#4}}
\newcommand{\eltlight}[2]{\ensuremath{\constr{exist}~#1~#2}}
\def\class{\kw{class}~}
\def\instance{\kw{instance}~}
\def\tclass#1{\coqdocind{#1}}
\def\where{\kw{where}~}
\def\module#1{\texttt{#1}}

\def\vector#1{\ind{vector}~#1}
\def\vnil{\constr{vnil}}
\def\vcons{\constr{vcons}}

\renewcommand{\bar}[1]{{\overline{#1}}}

\setlength{\coqdocbaseindent}{0em}

\renewcommand{\eqbr}{`=_{\beta}}

\def\coqlibrary#1#2#3{}

% Tutorial

\newboolean{darkfond}
\setboolean{darkfond}{false}
%\setboolean{darkfond}{true}

\definecolor{metavarcolor}{rgb}{0.5,0.0,1.0}
\definecolor{darkgreen}{rgb}{0.1,0.7,0.1}
\definecolor{answercolor}{rgb}{1.0,0.0,0.0}
\definecolor{normalcolor}{rgb}{0.0,0.0,0.0}
\definecolor{exbluecolor}{rgb}{0.1,0.1,0.9}
\definecolor{dontlookcolor}{rgb}{0.5,0.5,0.5}
\definecolor{termcolor}{rgb}{0.0,0.1,0.9}
\definecolor{lookcolor}{rgb}{0.8,0.2,0.0}
\definecolor{prooftermcolor}{rgb}{0.3,0.1,1.0}
\definecolor{typecolor}{rgb}{1.0,0.6,0.0}
\definecolor{taccolor}{rgb}{0.70,0.10,0.0}
\definecolor{pink}{rgb}{0.8,0.6,0.6}
\definecolor{darkmagenta}{rgb}{0.4,0.0,0.6}
\definecolor{darkblue}{rgb}{0.0,0.0,0.6}


\newcommand{\corkscrew}{\boldsymbol{\large \vdash}}
%\newcommand{\disj}{\mbox{$\backslash/$}}
%\newcommand{\conj}{\mbox{$/\backslash$}}
\newcommand{\juge}[3]{\mbox{$#1 \boldsymbol{\vdash} #2 : \color{red}#3 $}}
%\newcommand{\juge}[4]{\mbox{$#1,#2 \corkscrew #3 \boldsymbol{:} #4 $}}
%\newcommand{\smalljuge}[3]{\mbox{$#1 \corkscrew #2 \boldsymbol{:} #3 $}}
%\newcommand{\goal}[3]{\mbox{$#1,#2 \corkscrew^{\!\!\!?} #3  $}}
%\newcommand{\sgoal}[2]{\mbox{$#1\corkscrew^{\!\!\!\!?} #2 $}}
%\newcommand{\reduc}[5]{\mbox{$#1,#2 \corkscrew #3 \rhd_{#4}#5 $}}
%\newcommand{\convert}[5]{\mbox{$#1,#2 \corkscrew #3 =_{#4}#5 $}}
%\newcommand{\convorder}[5]{\mbox{$#1,#2 \corkscrew #3\leq _{#4}#5 $}}
%\newcommand{\wouff}[2]{\mbox{$\emph{WF}(#1)[#2]$}}


\newcommand{\mthese}{\underset{M}{\vdash}}
%\newcommand{\jthese}{\mathop{\vdash}_J}
%\newcommand{\jthese}{\stackrel[J]{}{\vdash}}
\newcommand{\jthese}{\underset{\tiny J}{\vdash}}
%\newcommand{\jthese}{\boldsymbol{\vdash}_{\tiny{ \!\!J}}}

\newcommand{\jsequent}[2]{\mbox{$#1\;\jthese\; #2 $}}
\newcommand{\msequent}[2]{\mbox{$#1\;\mthese\; #2 $}}



%\newcommand{\type}{\boldsymbol{:}}


\newcommand{\answ}{\color{answercolor}}
\newcommand{\coquser}{\color{normalcolor}}
\newcommand{\Ord}{\mbox{$\mathbb{O}$}}
\newcommand{\ords}{\mbox{\(\mathbb{O}\)}}
\newcommand{\nl}{\mbox{$\underset{\scriptstyle NL}{\longrightarrow}$}}
\newcommand{\base}[1]{{\color{lookcolor}#1}}
\newcommand{\om}{\mbox{$\omega$}}

%_{\alpha\beta\delta\iota}
\def\order{\mathcal{O}}
\def\tchecking#1#2#3{\ensuremath{#1 \vdash #2 : #3}}
\def\tconv#1#2#3#4{\ensuremath{#1 \vdash #2 = #3 "~>"
    #4}}
\def\tcumul#1#2#3#4{\ensuremath{#1 \vdash #2 \leq #3 "~>"
    #4}}

\def\tinfer#1#2#3#4#5#6{\ensuremath{#1; #2 \vdash #3 \Uparrow\ 
    "~>" #4 \vdash #5 : #6}}

\def\tcheck#1#2#3#4#5#6#7{\ensuremath{#1; #2 \vdash #3 \Downarrow #4
    "~>" #5 \vdash #6 : #7}}

\usepackage{mathpartir}
%\newcommand{\trule}[1]{\textsc{#1}}
%\newcommand{\irule}[3]
%{\inferrule{ #2 }{ #3 }{\ \trule{#1}}}
\newenvironment{subsecframe}[1]{\begin{subsection}{#1}\begin{frame}\frametitle{#1}}%
    {\end{frame}\end{subsection}}


\begin{document}


\setbeamertemplate{background canvas}[vertical shading][top=rouge1,middle=rouge1,bottom=rouge1]
\setbeamertemplate{footline}{\hspace{5em} \textcolor{white} {\null \hfill}\hspace{2em}\null \vspace*{3pt}}

\begin{frame}

\begin{textblock*}{40mm}[0,0](10mm,0mm)
 \includegraphics[width=6cm]{INRIA_CHERCHEURS_UK_RVB}
  \end{textblock*}

\begin{textblock*}{12cm}(13mm,50mm)
{\textcolor{white} {
{\huge \thetitle}\\[2mm]
{\theauthor}}}
\end{textblock*}


   \begin{textblock*}{40mm}[0,0](10mm,76mm)
  \begin{picture}(5,80)
\put(0,23){\includegraphics[width=4cm,height=1.5cm]{logobasrougeV1}}
\put(20,50){\tiny \textcolor{rouge2}{${}^{\dagger}$University of Pennsylvania}}
\put(20,40){\tiny \textcolor{rouge2}{${}^{*}$$πr^2$ - Inria \& PPS}}
\end{picture}
\end{textblock*}


\begin{textblock*}{7cm}(55mm,76mm)
{\textcolor{white}{{\theevent}}}
\end{textblock*}

\vspace*{-4pt}
\end{frame}


%%%%%%%%%%%%%%%%%%%%%%%%%%%

\setbeamertemplate{background canvas}[vertical shading][top=rouge1,bottom=rouge1,middle=rouge1]
\setbeamercolor{toto}{fg=blanc,bg=rouge1}

\setbeamertemplate{footline}
{
\begin{beamercolorbox}[wd=1\paperwidth,ht=15.5pt]{toto}
%\raisebox{1ex}{\includegraphics[width=25mm]{logobastrans}}
  \hspace{4em}
  \raisebox{2.5ex}
  {\thetitle{}  -- \theauthorabbr{}}\hfill 
  \raisebox{2.5ex}
  {\insertframenumber \hspace{2mm} \null }
\end{beamercolorbox}}

\setbeamertemplate{background canvas}[vertical shading][top=white,middle=white,bottom=white]

\def\bulletfail{\alert{\ding{54}}}
\def\bulletcheck{\ding{52}}

%\setbeamercolor{background}{fg=red,bg=white,text=blue}

%\setbeamercolor{subsection in toc shaded}{fg=gray,bg=white,text=blue}
\setbeamercolor{section in toc}{fg=black,bg=white,text=blue}
\setbeamercolor{subsection in toc}{fg=black,bg=white,text=blue}
%\setbeamertemplate{table of contents shaded}[default]
%\setbeamertemplate{subsection in toc shaded}{\textcolor{gray}}

\setbeamertemplate{section in toc shaded}
{\begin{colormixin}{20!white}{\usebeamertemplate{section in toc}}\end{colormixin}\unskip}

\setbeamertemplate{subsection in toc shaded}
{\begin{colormixin}{20!white}{\usebeamertemplate{subsection in toc}}\end{colormixin}\unskip}

\setbeamertemplate{subsubsection in toc shaded}
{\begin{colormixin}{20!white}{\usebeamertemplate{subsubsection in toc}}\end{colormixin}\unskip}


\frame<beamer>{\frametitle{\Coq 8.5}
\tableofcontents[sectionstyle=show/show,subsectionstyle=show/show/show]}

\section{Up and coming}
\begin{subsecframe}{Native compilation - M. Dénès}
\begin{itemize}
\item NbE style (?)
\item \Coq $"->"^{\texttt{nativecomp}}$ \textsc{ML} $"->"^{\texttt{ocamlopt}}$ \textsc{ASM}
\item Faster than \texttt{vm\_compute}
\end{itemize}

  \begin{itemize}
  \item[+] \emph{Optional}
  \item[++/--] \emph{Faster at runtime, compilation is slow}
  \end{itemize}
\end{subsecframe}

\begin{subsecframe}{Incremental development - E. Tassi}
  \begin{itemize}
  \item \texttt{git} model
  \item Granularity: opaque lemmas
  \item \texttt{.vi} interfaces: cross-file compilation
  \item CoqIDE only (emacs interface hard to implement)
  \end{itemize}

  \begin{itemize}
  \item[-/++] \emph{Not optional, backwards-compatible}
  \item[+++] \emph{Faster interaction and compilation}
  \end{itemize}
\end{subsecframe}

\begin{subsecframe}{Universe polymorphism - M. Sozeau}
% Def interactive
  \begin{itemize}
  \item \emph{Checked} universes.
  \item Global definitions or local with prenex, ``bounded''
    quantification. \emph{Conservativity proof}
  \item \emph{Elaboration} from typical ambiguity + implicit
    generalization.
  \item Checks \texttt{HoTT/Coq}, B. Barras and T. Coquand's semi-simplicial
    model and another (1-)groupoid model.
    \pause
  \item Better error messaging, unification sensitive to universes.
  \item All definitions and inductives can be polymorphic now.
  \end{itemize}

  \begin{itemize}
  \item[=/++] \emph{Optional except for checked
      universes, impacts the ML hacker only}. Backwards-compatibility layer.
  \item[=+/+] \emph{Comparable or better performance, more expressive}
  \end{itemize}
  
\end{subsecframe}

\begin{subsecframe}{Fast record projections - M. Sozeau}
  \begin{itemize}
  \item Economic representation
  \item Fast reduction
  \item No more stupid unfoldings
  \end{itemize}
  
  \begin{itemize}
  \item[+/+--] \emph{Optional, backwards-compatibility layer,
      small source-level incompatibilities}.
  \item[${+}^{ω}$] \emph{Exponentially better performance}
  \end{itemize}
  % * Graphe de perfs
\end{subsecframe}

\begin{subsecframe}{New proof engine - A. Spiwack}
  % * new internal proof engine (perfs)
  % - eauto backtrack with ; true proof search
  % "all:", tactical "+", "once", a native "refine"

  \begin{itemize}
  \item Existential variables + \texttt{LogicT}
  \item Proof-search semantics (\texttt{;}, \texttt{+}, $\ldots$)
  \item Cleaner code: abstract monadic interface
  \item Native \texttt{refine}, dependent goals, goal manipulation
  \end{itemize}
  
  \begin{itemize}
  \item[-/++] \emph{Not optional, backwards-compatibility layer}
  \item[--] \emph{0-15\% time overhead (still invesigating, \textsc{OCaml} issue likely)}
  \end{itemize}

\end{subsecframe}
\begin{subsecframe}{New coqdoc - Y. Régis-Gianas}
  \begin{itemize}
  \item Based on the AST and \alert{\emph{not}} lexing anymore
  \item All semantic information available
  \item Plugable interface
  \end{itemize}

  \begin{itemize}
  \item[+] \emph{Not optional but legacy \texttt{coqdoc} code included}
  \item[=] \emph{Same performance}
  \end{itemize}
  
  % \vfill
  % \begin{center}\alert{{\Huge Demo}}\end{center}
\end{subsecframe}
\end{document}

%%% Local Variables: 
%%% mode: latex
%%% TeX-PDF-mode: t
%%% TeX-master: "slides"
%%% End: 
